\hypertarget{cpython-memory-allocator-pymalloc-generational-garbage-collection}{%
\chapter{CPython Memory Allocator (PyMalloc) \& Generational Garbage
Collection}\label{cpython-memory-allocator-pymalloc-generational-garbage-collection}}

\hypertarget{cpythons-memory-allocation-engine-pymalloc}{%
\subsection{23.1 CPython's Memory Allocation Engine
(PyMalloc)}\label{cpythons-memory-allocation-engine-pymalloc}}

For large allocations (greater than 512 bytes), CPython forwards the
request directly to the system's standard C library allocator
(\passthrough{\lstinline!malloc()!}). However, for small objects
(\(\le 512\) bytes)which represent the vast majority of Python
allocationsstandard operating system allocators introduce high
fragmentation and locking overhead. To resolve this, CPython implements
a custom small-object allocator called \textbf{PyMalloc}.

\hypertarget{pymalloc-memory-hierarchy}{%
\subsubsection{1. PyMalloc Memory
Hierarchy}\label{pymalloc-memory-hierarchy}}

PyMalloc structures memory into three distinct layers to minimize
operating system allocation calls: * \textbf{Arenas (256 KB)}:
Contiguous memory blocks allocated from the operating system, aligned to
256 KB boundaries. Arenas manage memory at the virtual memory level and
contain exactly 64 pools. * \textbf{Pools (4 KB)}: Subdivisions of
arenas that match the operating system's virtual page size. Each pool is
dedicated to a single \textbf{size class} (e.g., all blocks in a pool
are 32 bytes, or all are 64 bytes). * \textbf{Blocks}: The actual chunks
of memory returned to the interpreter. Blocks range from 8 bytes to 512
bytes, aligned to 8-byte steps (giving 64 distinct size classes).

\hypertarget{c-level-struct-definitions-in-obmalloc.c}{%
\subsubsection{2. C-level Struct Definitions in
obmalloc.c}\label{c-level-struct-definitions-in-obmalloc.c}}

The structures for pools and arenas are defined in CPython's memory
management source file (\passthrough{\lstinline!Objects/obmalloc.c!}).

\hypertarget{pool-header-struct}{%
\paragraph{Pool Header Struct}\label{pool-header-struct}}

Every 4 KB pool begins with a header that tracks block allocations and
linked lists of free pools:

\begin{lstlisting}[language=C]
/* Objects/obmalloc.c */
typedef uint8_t block;

struct pool_header {
    union { 
        block *as_block; 
        uint32_t as_uint; 
    } nextfree;             /* Pointer to the next available free block in the pool */
    union { 
        block *as_block; 
        uint32_t as_uint; 
    } firstfree;            /* Pointer to the first free block in the pool */
    struct pool_header *nextpool; /* Link to the next pool in the active list */
    struct pool_header *prevpool; /* Link to the previous pool in the active list */
    uint32_t refcount;      /* Number of allocated blocks currently in use */
    uint32_t szidx;         /* Size class index of this pool */
    uint32_t freeblocks;    /* Number of free blocks remaining in the pool */
};
typedef struct pool_header *poolp;
\end{lstlisting}

\begin{itemize}
\tightlist
\item
  \passthrough{\lstinline!nextfree!}: Points to the next block in a
  singly linked list of free blocks inside the pool. PyMalloc uses
  \textbf{in-place pointers} inside the free blocks themselves to
  implement this list without consuming extra memory.
\item
  \passthrough{\lstinline!firstfree!}: Tracks the boundary of allocated
  blocks versus unallocated space in the pool.
\item
  \passthrough{\lstinline!szidx!}: Specifies which size class this pool
  belongs to.
\end{itemize}

\hypertarget{arena-object-struct}{%
\paragraph{Arena Object Struct}\label{arena-object-struct}}

CPython tracks arenas using an array of
\passthrough{\lstinline!arena\_object!} descriptors:

\begin{lstlisting}[language=C]
/* Objects/obmalloc.c */
struct arena_object {
    uintptr_t address;          /* Base virtual memory address of the 256 KB block */
    block* pool_address;        /* Pointer to the first pool inside this arena */
    uint32_t nfreepools;        /* Number of pools currently free in this arena */
    uint32_t ntotalpools;       /* Total number of pools inside this arena (usually 64) */
    struct arena_object* nextarena; /* Link to the next arena */
    struct arena_object* prevarena; /* Link to the previous arena */
};
\end{lstlisting}

\hypertarget{size-class-formula}{%
\subsubsection{3. Size Class Formula}\label{size-class-formula}}

The size class for an allocation request of size \(S\) is calculated by:

\[\text{Class Index} = \left\lceil \frac{S}{8} \right\rceil - 1\]

This mapped alignment restricts memory fragmentation, ensuring that
small object requests are answered in \(\mathcal{O}(1)\) time by
locating the active pool array matching the calculated class index.

\begin{center}\rule{0.5\linewidth}{0.5pt}\end{center}

\hypertarget{generational-garbage-collection-and-cycle-detection}{%
\subsection{23.2 Generational Garbage Collection and Cycle
Detection}\label{generational-garbage-collection-and-cycle-detection}}

While reference counting manages the vast majority of object lifecycles,
it cannot identify self-referencing cyclic loops (e.g.,
\passthrough{\lstinline!x = []; x.append(x)!}). To reclaim cyclic memory
leaks, CPython runs a cyclic Garbage Collector (GC) as a background
system.

\hypertarget{gc-memory-prefix-and-header-layout}{%
\subsubsection{1. GC Memory Prefix and Header
Layout}\label{gc-memory-prefix-and-header-layout}}

Every object tracked by the GC has an extra header prefix in memory
before its normal \passthrough{\lstinline!PyObject!} structure. When
\passthrough{\lstinline!PyObject\_GC\_New()!} is called, it allocates
memory block of size:

\[\text{Allocation Size} = \text{sizeof(PyGC\_Head)} + \text{sizeof(PyObject)}\]

\begin{lstlisting}
+--------------------------------------------------------+
|                      PyGC_Head                         |
|  - gc_next: Linkage to other GC tracked objects        |
|  - gc_prev: Linkage to other GC tracked objects        |
|  - gc_refs: Temporary reference state                  |
+--------------------------------------------------------+
|                      PyObject                          | <--- Object Pointer Returned to User
|  - ob_refcnt: Standard reference counter               |
|  - ob_type: Type pointer                               |
+--------------------------------------------------------+
\end{lstlisting}

The pointer returned to the user points directly to the
\passthrough{\lstinline!PyObject!} start. The VM accesses the GC header
by subtracting the header size:

\[\text{PyGC\_Head Pointer} = (\text{PyGC\_Head*})\text{op} - 1\]

The C-level layout of \passthrough{\lstinline!PyGC\_Head!} is defined in
\passthrough{\lstinline!Include/internal/pycore\_gc.h!}:

\begin{lstlisting}[language=C]
typedef union _gc_head {
    struct {
        uintptr_t gc_next;   /* Pointer to next object in generation list */
        uintptr_t gc_prev;   /* Pointer to previous object in generation list */
        Py_ssize_t gc_refs;  /* Reference count copy used during cycle checks */
    } gc;
    double dummy;            /* Forces double-word alignment adjustments */
} PyGC_Head;
\end{lstlisting}

\hypertarget{the-three-generations}{%
\subsubsection{2. The Three Generations}\label{the-three-generations}}

The GC divides tracked objects into three generations based on survival
age: * \textbf{Generation 0 (Gen 0)}: Where new objects are registered.
Collected frequently. * \textbf{Generation 1 (Gen 1)}: Objects that
survived a Gen 0 collection pass. * \textbf{Generation 2 (Gen 2)}:
Long-lived objects. Collected least frequently.

Collection is triggered when the number of allocations minus
deallocations in a generation exceeds a configured threshold:

\[\text{Allocations} - \text{Deallocations} > \text{Threshold}\]

Default thresholds are \passthrough{\lstinline!(700, 10, 10)!}. You can
view or change them via \passthrough{\lstinline!gc.get\_threshold()!}
and \passthrough{\lstinline!gc.set\_threshold()!}.

\begin{center}\rule{0.5\linewidth}{0.5pt}\end{center}

\hypertarget{the-cycle-detection-algorithm}{%
\subsection{23.3 The Cycle Detection
Algorithm}\label{the-cycle-detection-algorithm}}

To find and break reference cycles, the GC executes the following steps
during a collection pass:

\hypertarget{initializing-gc_refs}{%
\subsubsection{1. Initializing gc\_refs}\label{initializing-gc_refs}}

The GC copy-assigns the reference count of every object in the
collection generation to its \passthrough{\lstinline!gc\_refs!} field:

\[\text{gc\_refs} = \text{ob\_refcnt}\]

\hypertarget{traversing-internal-references}{%
\subsubsection{2. Traversing Internal
References}\label{traversing-internal-references}}

For each object in the generation list, the GC calls its
\passthrough{\lstinline!tp\_traverse!} slot (a C function that lists all
objects referenced by the object). If a referenced object is also in the
collection generation, the GC decrements that object's
\passthrough{\lstinline!gc\_refs!}:

\[\text{gc\_refs} = \text{gc\_refs} - 1\]

After traversing all objects, any object that still has
\passthrough{\lstinline!gc\_refs > 0!} must be reachable from outside
the generation list (e.g., from local variables, global parameters, or a
higher generation).

\hypertarget{resolving-and-splitting-collections}{%
\subsubsection{3. Resolving and Splitting
Collections}\label{resolving-and-splitting-collections}}

The GC splits the candidate list into two sets:
\passthrough{\lstinline!reachable!} and
\passthrough{\lstinline!unreachable!}.

\begin{lstlisting}
                  [GC Collection Candidates]
                              |
              Does object have gc_refs > 0?
                           /     \
                         Yes      No
                         /         \
            [Reachable Set]       [Unreachable Set]
\end{lstlisting}

\begin{enumerate}
\def\labelenumi{\arabic{enumi}.}
\tightlist
\item
  If an object has \passthrough{\lstinline!gc\_refs > 0!}, it is moved
  back to the \passthrough{\lstinline!reachable!} list.
\item
  All objects reachable \emph{from} that object (transitively traversed
  via \passthrough{\lstinline!tp\_traverse!}) are also marked as
  \passthrough{\lstinline!reachable!} and moved out of the
  \passthrough{\lstinline!unreachable!} list, even if their
  \passthrough{\lstinline!gc\_refs!} was zero.
\item
  Any remaining objects in the \passthrough{\lstinline!unreachable!} set
  are confirmed to be part of a cyclic garbage loop and are marked for
  deallocation.
\end{enumerate}

\hypertarget{clearing-cycles-and-deallocation}{%
\subsubsection{4. Clearing Cycles and
Deallocation}\label{clearing-cycles-and-deallocation}}

The GC breaks the reference cycle by calling the
\passthrough{\lstinline!tp\_clear!} slot on each object in the
\passthrough{\lstinline!unreachable!} set. This sets internal pointers
(like dictionary entries or list items) to
\passthrough{\lstinline!None!} or clears them, which decrements the
reference counts of the target objects and allows the standard reference
counter to deallocate the memory.

\begin{center}\rule{0.5\linewidth}{0.5pt}\end{center}
